% ID: SEL-D0032-OB-E03
% ESTADO: validada. S033, cotejada con examen y criterios.
Se toma como población del distribuidor la producción
indicada del mes. Como cada paquete pesa 1 kg, podemos
contar paquetes directamente a partir de los kilogramos.
La fábrica F1 aporta 8000 del tipo B y F2 aporta 10000.

\textbf{a)} Entre 45000 paquetes hay 18000 de tipo B.
Por tanto $P(B)=18000/45000=\boxed{0.40}$.

\textbf{b)} Si el paquete es A, restringimos el total a
$12000+15000=27000$. Las probabilidades condicionadas son
$12000/27000=4/9$ para F1 y $15000/27000=5/9$ para F2.
Es más probable que proceda de \textbf{F2}.
